\begin{frame}{Methodology}
    \huge
    \centering
    Methodology
\end{frame}

\begin{frame}{Methodology}
    HAHAHAHAHHA TODO because how on earth am I going to explain HJI derivation
    of optimal controls using the Lagrangian in conceivably 5 minutes?
    $n$ pursuers, $m$ evaders, something something something
\end{frame}

% In this paper, we consider a search-and-rescue mission
% in which $m$ ground agents are deployed to cover a large search area. We
% consider ground agents as independent and that they do not necessarily
% coordinate their movement with one another (mimicking a "divide-and-conquer"
% approach). During this mission, $n$ UAVs equipped with radio transmitters are
% tasked with forming an ad-hoc aerial network to maintain connectivity between
% all $m$ ground units. An optimal control input for the UAVs is desired to
% maximize connectivity and ensure that the ground units maintain the ability to
% communicate throughout the mission.

%
% We model this scenario as a two-player, zero-sum (2P0S), pursuit-evasion (PE)
% differential game in which UAVs are modelled as pursuers and ground agents as
% evaders. We follow the convention of modelling external system forces as an
% optimal adversary by considering the evaders as a competing adversary to the
% pursuers. We employ spectral graph theory to obtain a measure of network
% connectivity to be optimized, and then employ the Hamilton-Jacobi-Isaacs (HJI)
% partial differential equation (PDE) to determine a closed-form optimal control
% law for $n$ pursuers and $m$ evaders to maximize/minimize network connectivity
% respectively. We then construct a simulator to execute different PE scenarios,
% as well as a rendering engine for visually observing the behaviour of agents
% during each simulation run. We then visually and numerically evaluate the
% behaviour of the pursuers under their optimal control law when evaders follow a
% random walk.
